%------------------------------------------------------------------------------%
\begin{exercises}

% Exercício de prova do prof. Éden Amorim
%------------------------------------------------------------------------------%
\begin{exercise}
  Aplicando os testes de convergência verifique se a série numérica abaixo
  é convergente ou divergente
  \[
    \sum_{n=1}^{\infty} \frac{(-1)^{n+1}\; n}{3^n}
  \]
  \begin{answer}
    Pelo teste da razão
    \[
      \lim \abs{\frac{a_{n+1}}{a_n}} = \frac{1}{3} < 1
    \]
    portanto a série converge.
  \end{answer}
\end{exercise}

% Ex 1-8 pg 36 Thomas
%------------------------------------------------------------------------------%
\begin{exercise}
  Utilize o teste da razão para determinar se cada série converge ou diverge.
  \begin{mathlist}[2][series=ex-03-06-A]
    \item \sum_{n=1}^\infty \frac{2^n}{n!}
    \item \sum_{n=1}^\infty \frac{n+2}{3^n}
    \item \sum_{n=1}^\infty \frac{(n-1)!}{(n+1)^2}
    \item \sum_{n=1}^\infty \frac{2^{n+1}}{3^{n-1}\, n}
    \item \sum_{n=1}^\infty \frac{n^4}{4^n}
    \item \sum_{n=2}^\infty \frac{3^{n+2}}{\ln n}
    \item \sum_{n=1}^\infty \frac{n^2(n+2)!}{n!\,3^{2n}}
  \end{mathlist}
  \begin{mathlist}[1][resume=ex-03-06-A]
    \item \sum_{n=1}^\infty \frac{5^n\,n}{(2n+3)\ln(n+1)}
  \end{mathlist}
\end{exercise}

%------------------------------------------------------------------------------%
%------------------------------------------------------------------------------%
%------------------------------------------------------------------------------%
%------------------------------------------------------------------------------%
%------------------------------------------------------------------------------%
%------------------------------------------------------------------------------%

% Thomas pg 36
%------------------------------------------------------------------------------%
\begin{exercise}
  Determine se a série converge ou diverge.
\begin{mathlist}
  % 17-25
  \item \sum_{n=1}^\infty \frac{n^{\sqrt{2}}}{2^n}
  \item \sum_{n=1}^\infty n^2 e^{-n}
  \item \sum_{n=1}^\infty n! e^{-n}
  \item \sum_{n=1}^\infty \frac{n!}{10^n}
  \item \sum_{n=1}^\infty \frac{n^{10}}{10^n}
  \item \sum_{n=1}^\infty \left(\frac{n-2}{n}\right)^n
  \item \sum_{n=1}^\infty \frac{2+(-1)^n}{1,25^n}
  \item \sum_{n=1}^\infty \frac{(-2)^n}{3^n}
  \item \sum_{n=1}^\infty \left(1-\frac{3}{n}\right)^n
  % 26-34
  \item \sum_{n=1}^\infty \left(1-\frac{1}{3n}\right)^n
  \item \sum_{n=1}^\infty \frac{n\ln(n)}{2^n}
  \item \sum_{n=1}^\infty \frac{(n+1)(n+2)}{n!}
  \item \sum_{n=1}^\infty e^{-n}\,n^3
  % 35-39
  \item \sum_{n=1}^\infty \frac{(n+3)!}{3!\,n!\,3^n}
  \item \sum_{n=1}^\infty \frac{n\,2^n\,(n+1)!}{3^n\,n!}
  \item \sum_{n=1}^\infty \frac{n!}{(2n+1)!}
  \item \sum_{n=1}^\infty \frac{n!}{n^n}
  % 40-44
  \item \sum_{n=1}^\infty \frac{3^n}{n^32^n}
  \item \sum_{n=1}^\infty \frac{(n!)^2}{(2n)!}
  \item \sum_{n=1}^\infty \frac{(2n\!+\!3)(2^n\!+\!3)}{3^n+2}
\end{mathlist}
\end{exercise}

% Thomas pg 36
%------------------------------------------------------------------------------%
\begin{exercise}
Determine se as séries
\(
   \sum_{n=1}^\infty a_n
\)
cujos termos são definidos pelas relações de recorrência
são convergentes ou divergentes.
\begin{mathlist}
  % 45 - 54
\item a_1 = 2           \) \tabto{30mm} \( a_{n+1} = \frac{1+\sin(n)}{n} a_n    \)
\item a_1 = 1           \) \tabto{30mm} \( a_{n+1} = \frac{1+\arctan(n)}{n} a_n \)
\item a_1 = \frac{1}{3} \) \tabto{30mm} \( a_{n+1} = \frac{3n-1}{2n+5} a_n      \)
\item a_1 = 3           \) \tabto{30mm} \( a_{n+1} = \frac{n}{n+1} a_n          \)
\item a_1 = 2           \) \tabto{30mm} \( a_{n+1} = \frac{2}{n} a_n            \)
\item a_1 = 5           \) \tabto{30mm} \( a_{n+1} = \frac{\sqrt[n]{n}}{2} a_n  \)
\item a_1 = 1           \) \tabto{30mm} \( a_{n+1} = \frac{1+\ln(n)}{n} a_n     \)
\item a_1 = \frac{1}{2} \) \tabto{30mm} \( a_{n+1} = \frac{n+\ln(n)}{n+10} a_n  \)
\item a_1 = \frac{1}{3} \) \tabto{30mm} \( a_{n+1} = \sqrt[n]{a_n}              \)
\item a_1 = \frac{1}{2} \) \tabto{30mm} \( a_{n+1} = \left(a_n\right)^{n+1}     \)
\end{mathlist}
\end{exercise}

%------------------------------------------------------------------------------%
\end{exercises}
%------------------------------------------------------------------------------%
